\documentclass[11pt]{article}
\usepackage[letterpaper,margin=1in]{geometry}
\usepackage{amsmath,amssymb,amsthm,mathtools}
\usepackage{booktabs,array,longtable}
\usepackage{xcolor}
\usepackage{microtype}
\usepackage[hidelinks]{hyperref}
\hypersetup{
  pdftitle={An ultraviolet overlap benchmark for a nonlinear SU(2) lattice block},
  pdfsubject={Heat-kernel tails, Haar Jacobians, and restricted polymer activities}
}

\newtheorem{theorem}{Theorem}[section]
\newtheorem{proposition}[theorem]{Proposition}
\newtheorem{lemma}[theorem]{Lemma}
\newtheorem{corollary}[theorem]{Corollary}
\theoremstyle{definition}
\newtheorem{definition}[theorem]{Definition}
\newtheorem{remark}[theorem]{Remark}

\newcommand{\SU}{\mathrm{SU}(2)}
\newcommand{\R}{\mathbb{R}}
\newcommand{\norm}[1]{\left\lVert #1\right\rVert}
\newcommand{\dd}{\,\mathrm{d}}
\newcommand{\proved}{\textcolor{green!45!black}{\textsf{proved}}}
\newcommand{\restricted}{\textcolor{blue!65!black}{\textsf{restricted}}}
\newcommand{\openstatus}{\textcolor{red!70!black}{\textsf{open}}}

\title{An Ultraviolet Overlap Benchmark for a Nonlinear $\SU$ Lattice Block\\
\large Heat-kernel tails, Haar Jacobians, and the first dependent-measure obstruction}
\author{Mathematical research dossier}
\date{20 July 2026}

\begin{document}
\maketitle

\begin{abstract}
The preceding nonlinear block certificate controlled a very small local
chart, but at heat time $t=10^{-2}$ almost every standalone heat-kernel
plaquette lay outside it.  We determine whether that mismatch is structural
or merely a scale mismatch.  First, a cancellation-free Poisson-image
estimate bounds the standalone $\SU$ curvature tail uniformly at weak
coupling.  Second, Euler's product for $\sin r/r$ gives explicit Hessian and
third-derivative bounds for the link Haar Jacobian.  Third, an independent
connected-set theorem converts a one-plaquette tail into a weighted block
activity benchmark.

At
\[
 t=3\cdot10^{-6},\qquad r=0.0041,\qquad \rho=2r=0.0082,
 \qquad r_F=10^{-5},
\]
all local-action inequalities remain strict after the Haar contribution.  The
standalone tail is at most $5.51231\cdot10^{-5}$; after a 216-plaquette union
bound, the weighted independent-block branching factor is at most $0.842163$
and the rooted weighted activity sum is at most $0.083370$.  Thus local
nonlinear control and small ultraviolet defect activities are not
intrinsically incompatible.

This is a restricted coexistence certificate, not full-measure Yang--Mills
control.  Plaquettes in the four-dimensional link measure share links, obey
nonabelian constraints, and carry boundary and topological information.  The
standalone heat-kernel law is not proved here to be their marginal, and block
independence is false.  The first unproved implication is therefore a
volume-uniform dependent activity theorem for the complete gauge-invariant
link measure.
\end{abstract}

\tableofcontents

\section{Question, conventions, and result}

We use the unit-three-sphere model of $\SU$ and heat generator
$\tfrac12\Delta$.  A class angle is $\theta\in[0,\pi]$, and the leading
small-time action is $\theta^2/(2t)$.  The side-two four-dimensional cubical
block has 216 links and 216 plaquettes.  Stage 08 established the nonlinear
plaquette-log and local Coulomb-chart constants on this block; its exact
quaternion calculus is the finite geometric input used below.

The distinction between three measures is essential.
\begin{enumerate}
\item The \emph{standalone heat-kernel law} has class density
  $p_t(\theta)(2/\pi)\sin^2\theta\dd\theta$.
\item The \emph{independent-block benchmark} assigns independent bad events
  with a declared block-probability upper bound.
\item The \emph{four-dimensional link measure} integrates link variables
  against a product of plaquette weights and link Haar measures.  Its
  plaquette variables are constrained and dependent.
\end{enumerate}
Only the first two enter the probability theorem proved here.  The local
action inequalities, by contrast, concern the actual finite block action as
a function of its link coordinates.

\begin{theorem}[Restricted ultraviolet overlap witness]\label{thm:main}
Let
\[
 t=3\cdot10^{-6},\quad r=0.0041,\quad \rho=0.0082,
 \quad r_F=10^{-5},\quad \Delta=8,\quad \alpha=0.1.
\]
Then the following statements hold.
\begin{enumerate}
\item On the Stage 08 logarithmic Coulomb chart, the nonlinear physical
  singular-value lower bound is $1.332871$.
\item After adding the exponential-link Haar density, the scaled Hessian
  remainder is at most $0.968684<1$, the centered fiber transport coefficient
  is at most $0.003941<1$, and the generated reference-base margin is at
  least $0.060732>0$.
\item For a standalone heat-kernel plaquette,
  \[
    \Pr_t\{\theta\ge\rho\}\le 5.512301\cdot10^{-5}.
  \]
\item Let $p_B$ be the union-bound proxy for 216 such plaquette failures.
  Then $p_B\le0.0119066$.  For independent bad blocks on a graph of maximum
  degree eight,
  \[
   e^\alpha\Delta^2p_B\le0.842163<1
  \]
  and
  \[
   \sum_{X\ni o}e^{\alpha|X|}\Pr\{X\text{ is all bad}\}
   \le 0.083370<1.
  \]
\end{enumerate}
Every displayed pass inequality is strict; hence their simultaneous validity
persists on a nonempty open neighborhood of $(t,r,r_F)$.
\end{theorem}

The proof is assembled in Sections~\ref{sec:haar}--\ref{sec:overlap}.

\begin{remark}[Logical status]
Part 4 is an independent-block theorem.  It is not applied to the dependent
four-dimensional link measure.  Theorem~\ref{thm:main} proves numerical and
analytic coexistence of the two required regimes, not basin entry for the
complete theory.
\end{remark}

\section{The link Haar Jacobian}\label{sec:haar}

For $A\in\R^3$ with $a=\norm A<\pi$, exponential coordinates pull normalized
Haar measure back, up to a constant independent of $A$, with radial Jacobian
\[
 J(A)=\left(\frac{\sin a}{a}\right)^2.
\]
Its contribution to the coordinate action is
\[
 h(A)=-\log J(A)=-2\log\frac{\sin a}{a}.
\]

\begin{lemma}[Local Haar derivative bounds]\label{lem:haar}
For $0\le a\le r<\pi$, put $x=r/\pi$ and $q=x^2$.  Then
\begin{align}
 \norm{D^2h(A)}
 &\le H_2(r):=\frac23\frac{1+q}{(1-q)^2},\label{eq:H2}\\
 \norm{D^3h(A)}
 &\le H_3(r):=\frac{2x}{3\pi}
 \left(\frac{4(1+q)}{(1-q)^3}+\frac{2}{(1-q)^2}\right).
 \label{eq:H3}
\end{align}
At $r=0.0041$ these bounds are $H_2\le0.666671$ and
$H_3\le0.001662$.
\end{lemma}

\begin{proof}
Euler's product and absolute convergence on compact subsets of $(-\pi,\pi)$
give
\[
 h(a)=2\sum_{n\ge1}\frac{\zeta(2n)}{n\pi^{2n}}a^{2n}.
\]
For $p=2n$, the Hessian operator norm of $\norm A^p$ is at most
$p(p-1)a^{p-2}$.  Direct differentiation gives the third tensor as one radial
cubic term and three mixed terms; the triangle inequality bounds its operator
norm by $p(p-1)(p-2)a^{p-3}$.  Since
$\zeta(2n)\le\zeta(2)=\pi^2/6$, termwise differentiation yields
\[
 \norm{D^2h}
 \le \frac23\sum_{n\ge1}(2n-1)q^{n-1}
 =\frac23\frac{1+q}{(1-q)^2}.
\]
Likewise, after writing $m=n-1$,
\[
 \norm{D^3h}
 \le\frac{2}{3\pi x}
 \sum_{m\ge1}(2m+1)(2m)q^m,
\]
and the standard sums for $\sum m q^m$ and $\sum m^2q^m$ give
\eqref{eq:H3}.
\end{proof}

The scaled action used in the local norm is $tS$.  Since the link Haar
potential is a sum of one-link functions, its global Hessian and third-tensor
bounds are $tH_2(r)$ and $tH_3(r)$; no plaquette-count factor is introduced.
At the witness these contributions are $2.00002\cdot10^{-6}$ and less than
$5\cdot10^{-9}$, respectively.

\section{A cancellation-free standalone tail}\label{sec:tail}

The character expansion is exact but becomes a numerically ill-conditioned
alternating representation when $t$ is extremely small.  The Poisson-image
representation exposes the ultraviolet concentration directly.  Up to a
positive factor independent of $\theta$, the heat-kernel shape is
\begin{equation}\label{eq:image}
 p_t(\theta)\propto
 \frac1{\sin\theta}\sum_{k\in\mathbb Z}
 (\theta+2\pi k)
 \exp\left[-\frac{(\theta+2\pi k)^2}{2t}\right].
\end{equation}
Together with radial Haar density, the unnormalized radial mass is therefore
$\sin\theta$ times the image sum.

\begin{proposition}[Image tail upper bound]\label{prop:tail}
Let $0<\rho<\pi$, $t>0$, and
$L=\min\{\pi/2,3\sqrt t\}$.  Define
\begin{align*}
 D_0(t,L)&=\frac{\sin L}{L}
 \left[\sqrt{\frac\pi2}t^{3/2}
 \operatorname{erf}\left(\frac{L}{\sqrt{2t}}\right)
 -tL e^{-L^2/(2t)}\right],\\
 N_0(t,\rho)&=t\rho e^{-\rho^2/(2t)}
 +t\sqrt{\frac{\pi t}{2}}
 \operatorname{erfc}\left(\frac{\rho}{\sqrt{2t}}\right),\\
 E(t)&=t\frac{e^{-\pi^2/(2t)}+e^{-2\pi^2/t}}
 {1-e^{-4\pi^2/t}}.
\end{align*}
If $D_0(t,L)>E(t)$, then the standalone heat-kernel law satisfies
\begin{equation}\label{eq:tailbound}
 \Pr_t\{\theta\ge\rho\}
 \le \min\left\{1,\frac{N_0(t,\rho)+E(t)}{D_0(t,L)-E(t)}\right\}.
\end{equation}
\end{proposition}

\begin{proof}
For the zero image, $\sin\theta\le\theta$ and extension of the upper
integration limit to infinity give
\[
 \int_\rho^\pi \theta\sin\theta e^{-\theta^2/(2t)}\dd\theta
 \le\int_\rho^\infty\theta^2e^{-\theta^2/(2t)}\dd\theta=N_0(t,\rho).
\]
On $[0,L]$, concavity of sine gives
$\sin\theta\ge(\sin L/L)\theta$, so the retained zero-image normalization is
at least $D_0(t,L)$.

It remains to bound every nonzero image in absolute value.  For the negative
$j$th image, $j\ge1$, substitution gives an integral bounded by
$t e^{-(2j-1)^2\pi^2/(2t)}$; for the positive image the bound is
$t e^{-2j^2\pi^2/t}$.  Consecutive terms in both sequences have ratio at most
$e^{-4\pi^2/t}$.  Summing the geometric majorant gives $E(t)$.  Adding $E$ to
the tail numerator and subtracting it from the normalization denominator
proves \eqref{eq:tailbound}.
\end{proof}

At the witness, $\rho/\sqrt t=4.734272$ and the nonzero-image term is smaller
than $10^{-714000}$.  Formula~\eqref{eq:tailbound} gives
$5.512301\cdot10^{-5}$.  An independent character calculation truncated at
representation index 6000 encloses the actual tail between
$5.3507587609\cdot10^{-5}$ and $5.3507587704\cdot10^{-5}$, providing a
cross-representation check without entering the proof of the upper bound.

\section{The restricted weighted-polymer theorem}\label{sec:polymer}

\begin{theorem}[Independent rooted connected sets]\label{thm:polymer}
Let a graph have maximum degree $\Delta$.  Suppose vertices are bad
independently with probability at most $p_B$.  For $\alpha\ge0$, if
$q_\alpha=e^\alpha\Delta^2p_B<1$, then
\begin{equation}\label{eq:polymer}
 \sum_{\substack{X\ni o\\X\text{ connected}}}
 e^{\alpha|X|}\Pr\{X\text{ is all bad}\}
 \le \frac{e^\alpha p_B}{1-q_\alpha}.
\end{equation}
\end{theorem}

\begin{proof}
A depth-first traversal encodes a rooted connected $n$-set by a walk of
length at most $2(n-1)$, so the number of such sets is at most
$\Delta^{2(n-1)}$.  Independence bounds the all-bad probability by $p_B^n$.
Summing the resulting geometric series proves \eqref{eq:polymer}.
\end{proof}

For a block containing 216 plaquettes, the union bound gives
$p_B\le216p$.  No within-block independence is needed for this union bound.
Independence between bad blocks is, however, indispensable for
Theorem~\ref{thm:polymer}; it is the assumption that fails in the target
four-dimensional gauge measure.

\section{Composition and numerical margins}\label{sec:overlap}

Let $\eta_P$ and $\tau_P$ be the Stage 08 scaled plaquette-action remainder
and third-tensor bound.  Lemma~\ref{lem:haar} gives the augmented quantities
\[
 \eta=\eta_P+tH_2(r),\qquad \tau=\tau_P+tH_3(r).
\]
In the same Hodge-whitened retained/fiber coordinates, with reference
curvatures $\lambda_F=2$, $\lambda_Y=1$, the five-component bounds become
\[
 \varepsilon_F=\eta,\quad
 \varepsilon_{FY}=\frac{\eta}{\sqrt2},\quad
 \varepsilon_Y=\frac\eta2,\quad
 \tau_F=\tau,\quad \tau_{FY}=\frac\tau{\sqrt2}.
\]
Writing $m_F=2-\eta$, the response and generated-action formulas are
\begin{align*}
 K_F&=\frac{r_F}{m_F}
 \left(\frac{\tau}{\sqrt2}
 +\tau\frac{\eta/\sqrt2}{m_F}\right),\\
 \eta_{\mathrm{gen}}&=\frac\eta2+
 \frac{(\eta/\sqrt2)^2}{m_F},
 \qquad m_Y=1-\eta_{\mathrm{gen}}.
\end{align*}
Substitution gives the values in Table~\ref{tab:witness}.

\begin{table}[ht]
\centering
\caption{Certified witness values.}\label{tab:witness}
\begin{tabular}{@{}lll@{}}
\toprule
Quantity & Bound & Test \\
\midrule
$\eta_P$ & $0.968681008$ & local plaquette input \\
$tH_2(r)$ & $0.000002001$ & Haar contribution \\
$\eta$ & $0.968683008$ & $<1$ \\
$m_F$ & $1.031316992$ & $>0$ \\
$K_F$ & $0.003940979$ & $<1$ \\
$m_Y$ & $0.060732040$ & $>0$ \\
$p$ & $0.0000551231$ & standalone tail \\
$p_B$ & $0.011906570$ & $216p$ \\
$e^\alpha\Delta^2p_B$ & $0.842162811$ & $<1$ \\
rooted weighted sum & $0.083369414$ & $<1$ \\
\bottomrule
\end{tabular}
\end{table}

This proves Theorem~\ref{thm:main}.  The strict-continuity conclusion follows
because the denominators in the local inverse, fiber response, image
normalization, and polymer sum remain positive at the witness.

For comparison, at the Stage 08 point $(t,r)=(10^{-2},0.003)$, the exact
standalone probability outside $\rho=0.006$ exceeds $0.9999$.  The local
action remains controlled, but the independent activity proxy fails
completely.  The successful ultraviolet witness shows that this failure was
not a universal chart-versus-probability incompatibility.

\section{What the computation does and does not certify}

The executable checker evaluates the exact displayed formulas, performs a
declared $25\times25$ parameter-grid calibration, and runs positive,
boundary, monotonicity, cross-representation, and deliberately failing tests.
Of 625 grid points in
\[
 t\in[10^{-6},5\cdot10^{-6}],\qquad
 r\in[0.0036,0.0042],
\]
389 pass the restricted inequalities with $r_F=10^{-5}$.  This demonstrates
that the witness is not an isolated floating-point accident.  The grid search
does not claim global optimality, and no proof depends on the selected grid.

The mechanization boundary is equally important.  The checker does not
formalize the analytic Poisson summation identity, the Stage 08 nonlinear
inverse theorem, or any infinite-volume cluster expansion.  It checks their
closed finite consequences and fails closed on invalid radii, nonpositive
normalizations, exhausted coercivity, and divergent polymer denominators.

\section{First unproved implication and strategic consequence}

The target is now a dependent chart--polymer overlap theorem.  For the actual
regulated link measure $\mu_{t,L}$ and a gauge-invariant exceptional event,
one needs activities $z_{t,L}(X)$ satisfying, in the same local-control
window,
\begin{equation}\label{eq:target}
 \sup_L\sup_{o}\sum_{X\ni o}
 e^{\alpha|X|}|z_{t,L}(X)|<1.
\end{equation}
The estimate must retain the following data rather than absorbing them into a
generic ``large-field error'':
\begin{itemize}
\item shared-link dependence and nonabelian Bianchi constraints;
\item the orbit-space atlas and any Coulomb-slice/Faddeev--Popov density;
\item conditional boundary holonomy;
\item center flux, torons, instanton-like sectors, and global topology;
\item constants uniform in volume, boundary condition, and the repeated RG
  scale.
\end{itemize}

The strategic change is sharp.  Enlarging the local chart is no longer the
only apparent route.  The ultraviolet scale already separates the typical
curvature, of order $\sqrt t$, from a chart radius that remains small enough
for nonlinear coercivity.  The next work should exploit the actual
gauge-covariant dependence structure to prove \eqref{eq:target}, or else show
precisely why the chosen block event and norm fail under conditioning.

\section{Frontier audit}

\begin{longtable}{@{}p{0.47\textwidth}p{0.16\textwidth}p{0.27\textwidth}@{}}
\toprule
Claim & Status & Scope \\
\midrule
\endhead
Poisson-image standalone tail bound & \proved & one $\SU$ heat-kernel law \\
Link Haar derivative bounds & \proved & local exponential chart \\
Independent rooted weighted-polymer sum & \proved & independent block events \\
Simultaneous strict witness & \restricted & coexistence benchmark \\
Tail for the full 4D plaquette marginal & \openstatus & interacting link measure \\
Dependent activity bound \eqref{eq:target} & \openstatus & volume-uniform 4D measure \\
Repeated RG to a finite physical scale & \openstatus & cutoff- and scale-uniform \\
Continuum construction, OS reconstruction, and gap & \openstatus & Clay endpoint \\
\bottomrule
\end{longtable}

The strongest defensible level is \textbf{C2, restricted theorem}.  This
document does not construct four-dimensional continuum Yang--Mills theory and
does not prove a Hamiltonian mass gap.  The existence and mass-gap problem
remains open.

\begin{thebibliography}{9}
\bibitem{MenottiOnofri1981}
P.~Menotti and E.~Onofri,
``The action of $SU(N)$ lattice gauge theory in terms of the heat kernel on the
group manifold,''
\emph{Nuclear Physics B} \textbf{190} (1981), 288--300.

\bibitem{Balaban1984}
T.~Balaban,
``Propagators and renormalization transformations for lattice gauge
theories. I,''
\emph{Communications in Mathematical Physics} \textbf{95} (1984), 17--40.

\bibitem{Balaban1987}
T.~Balaban,
``Renormalization group approach to lattice gauge field theories. I.
Generation of effective actions in a small field approximation and a coupling
constant renormalization in four dimensions,''
\emph{Communications in Mathematical Physics} \textbf{109} (1987), 249--301.
\end{thebibliography}

\end{document}
